\input{estilo}
% Mejorada, esquemática: un dato entra al mapeo y sale la máquina donde vive;
% el mapeo puede ser una tabla (una grilla que alguien mantiene, como en GFS)
% o una función (que se calcula, como en hash consistente), simétricas bajo
% el mapeo.
\begin{document}
\begin{figura}\begin{tikzpicture}[maq/.style={caja, minimum width=7mm, minimum height=6mm}]
  \node[caja, minimum width=14mm] (d) at (0,0) {dato};
  \node[caja, minimum width=16mm, minimum height=9mm] (m) at (3.0,0) {mapeo};
  \draw[msg] (d) -- (m);
  % la ubicación física: una de varias máquinas
  \foreach \i/\y in {1/0.8, 2/0, 3/-0.8} \node[maq] (q\i) at (6.0,\y) {};
  \node[maq, fill=segundo!25, draw=segundo] at (6.0,0) {};
  \node[etiq, align=center, right=4pt] at (q2.east) {ubicación\\física};
  \draw[msg, draw=segundo] (m) -- (q2);
  % las dos maneras
  % la grilla, con cada línea trazada una sola vez para que todas pesen igual
  \draw[fill=white] (1.5,-2.55) rectangle (2.5,-1.65);
  \draw (2.0,-2.55) -- (2.0,-1.65) (1.5,-1.95) -- (2.5,-1.95) (1.5,-2.25) -- (2.5,-2.25);
  \coordinate (tb) at (2.0,-1.6);
  \node[etiq, anchor=base] (tl) at (2.0,-2.95) {tabla};
  \node[caja, minimum width=10mm, minimum height=9mm] (f) at (4.0,-2.1) {\var{f}};
  \node[etiq, anchor=base] (fl) at (4.0,-2.95) {función};
  \draw[respuesta] (m.south) -- (tb);
  \draw[respuesta] (m.south) -- (f.north);
\end{tikzpicture}\end{figura}
\end{document}
